\subsection{Simply connected 4-manifolds and stabilization}
\label{subsec:low-dimensional-topology}

The calculation of the invariant will use the following classical results
only through the precise interfaces recorded here.

\begin{definition}[Intersection lattice and hyperbolic stabilization]
\label{def:intersection-lattice-stabilization}
For a closed, connected, oriented, simply connected 4-manifold $X$, set
\[
  L_X:=\bigl(H_2(X;\mathbb Z),Q_X\bigr),
\]
where $Q_X$ is the integral intersection pairing.  Let
\[
  H:=L_{S^2\times S^2}
   =\left(\mathbb Z^2,
      \begin{pmatrix}0&1\\1&0\end{pmatrix}\right),
  \qquad
  X^{(r)}:=X\mathbin{\#}r(S^2\times S^2).
\]
\end{definition}

\begin{lemma}[Intersection lattice under stabilization]
\label{lem:intersection-lattice-stabilization}
Let \(X\) be a closed, connected, oriented, simply connected
4-manifold.  For every \(r\geq0\),
\[
  L_{X^{(r)}}\cong L_X\oplus rH.
\]
\end{lemma}

\begin{proof}
The second homology and intersection pairing of an oriented connected sum are
the orthogonal direct sums of those of its summands.
\end{proof}

\begin{remark}[Intersection-form conventions]
\label{rem:intersection-form-conventions}
The intersection pairing and connected-sum convention are those of
\cite[Definition~1.2.1 and Section~1.3]{GompfStipsicz1999}.
\end{remark}

\begin{definition}[Marked intersection lattice]
\label{def:marked-intersection-lattice}
For a simply connected marked 4-manifold $(X,\alpha)$, its
\emph{marked intersection lattice} is
\[
  \bigl(H_2(X;\mathbb Z),Q_X,\alpha\bigr).
\]
An isomorphism of marked intersection lattices means an isometry of the
intersection forms that carries one distinguished class to the other.
\end{definition}

\begin{remark}[Pointed lattice convention]
\label{rem:pointed-lattice-convention}
This is the pointed version of the lattice in
\cite[Definition~1.2.1]{GompfStipsicz1999}.
\end{remark}

\begin{theorem}[Classification of indefinite unimodular forms]
\label{thm:indefinite-unimodular-forms}
Let $L$ be an indefinite unimodular integral lattice.
\begin{enumerate}[label=(\roman*),ref=\roman*]
  \item If $L$ is odd, then
  \[
    L\cong b_+(L)\langle 1\rangle
      \oplus b_-(L)\langle-1\rangle.
  \]
  \item If $L$ is even, then $L$ is a direct sum of copies of the
  hyperbolic plane $H$ and copies of $E_8$ or $-E_8$, with the numbers
  determined by its rank and signature.
\end{enumerate}
\end{theorem}

\begin{proof}
This is \cite[Theorem~1.2.21]{GompfStipsicz1999}; see also
\cite[Theorem~2.3.1]{PetersSterk2024}.
\end{proof}

\begin{theorem}[Freedman and Donaldson]
\label{thm:freedman-donaldson}
Let \(X\) and \(Y\) be closed, oriented, simply connected smooth
4-manifolds.
\begin{enumerate}[label=(\roman*),ref=\roman*]
  \item Every isometry \(L_X\to L_Y\) is induced by an
  orientation-preserving homeomorphism \(X\to Y\).
  \item If the intersection form of \(X\) is definite, then it is diagonalizable
  over $\mathbb Z$.
\end{enumerate}
\end{theorem}

\begin{proof}
The first statement is the smoothable case of Freedman's classification
\cite[Theorem~1.2.27]{GompfStipsicz1999,Freedman1982}.  The second is
Donaldson's diagonalization theorem
\cite[Theorem~1.2.30]{GompfStipsicz1999,Donaldson1983}.
\end{proof}

\begin{theorem}[Wall stable classification]
\label{thm:wall-stable-classification}
Let $X$ and $Y$ be closed, oriented, simply connected smooth
4-manifolds.  If $L_X\cong L_Y$, then there is an integer $r\geq 0$ and
an orientation-preserving diffeomorphism
\[
  X^{(r)}\cong Y^{(r)}.
\]
\end{theorem}

\begin{proof}
Wall proves that an isometry of the intersection forms gives an
$h$-cobordism \cite[Theorem~2]{Wall1964}, and that $h$-cobordant simply
connected 4-manifolds become diffeomorphic after taking connected sums
with the same number of copies of $S^2\times S^2$
\cite[Theorem~3]{Wall1964}.
\end{proof}

\begin{theorem}[Ruberman--Strle one-stable realization]
\label{thm:ruberman-strle-realization}
Let $X$ be a closed, oriented, simply connected smooth 4-manifold and let
$A\in O(Q_X)$ be an automorphism of its intersection form.  Then
\[
  A\oplus\operatorname{id}_H
\]
is induced by an orientation-preserving self-diffeomorphism of
$X\#(S^2\times S^2)$.
\end{theorem}

\begin{proof}
This is \cite[Theorem~2.1]{RubermanStrle2022}.
\end{proof}

\begin{corollary}[Marked stable realization]
\label{cor:marked-stable-realization}
Let $f\colon X\to Y$ be an orientation-preserving homeomorphism between
closed, oriented, simply connected smooth 4-manifolds, and let
$\alpha\in H_2(X;\mathbb Z)$.  For some $r\geq 0$, there is an
orientation-preserving diffeomorphism
\[
  F\colon X^{(r)}\longrightarrow Y^{(r)}
\]
such that
\[
  F_*(\alpha\oplus0)=f_*(\alpha)\oplus0.
\]
\end{corollary}

\begin{proof}
By \cref{thm:wall-stable-classification}, choose
$G\colon X^{(k)}\to Y^{(k)}$.  The isometry
\[
  B:=(f_*\oplus\operatorname{id}_{kH})G_*^{-1}
  \in O(Q_{Y^{(k)}})
\]
is realized after one further stabilization by
\cref{thm:ruberman-strle-realization}.  Composing that realization with
$G\#\operatorname{id}_{S^2\times S^2}$ gives the required $F$ with
$r=k+1$.
\end{proof}

\begin{theorem}[Reutter's stable-evaluation principle]
\label{thm:reutter-stable-evaluation}
An indecomposable semisimple oriented four-dimensional TQFT assigns an
invertible scalar to $S^2\times S^2$ and takes the same value on stably
diffeomorphic closed 4-manifolds.  Consequently, its values on simply
connected 4-manifolds reduce to its values on standard connected-sum
blocks.
\end{theorem}

\begin{proof}
Stable-diffeomorphism invariance is \cite[Theorem~A]{Reutter2023},
invertibility of the $S^2\times S^2$ value is
\cite[Theorem~3.5]{Reutter2023}, and the resulting simply connected formulas
are recorded in \cite[Corollary~3.9]{Reutter2023}.
\end{proof}

\begin{remark}[Use in the twisted theory]
\label{rem:twisted-stabilization-use}
\Cref{cor:marked-stable-realization} supplies the marked topological input
needed later.  Reutter's theorem is cited as the untwisted precedent for the
stabilization strategy; it does not by itself prove any statement about the
homologically twisted evaluator.
\end{remark}
